% !TEX program = pdflatex
\documentclass[12pt,a4paper]{article}

% ── Font & tiếng Việt (pdfLaTeX) ──────────────────────────────────────────────
\usepackage[utf8]{inputenc}
\usepackage[T5]{fontenc}
\usepackage[vietnamese]{babel}
\usepackage{mathptmx} % Font Times (tương tự Times New Roman)
\usepackage{helvet}   % Font Arial/Helvetica cho sans-serif
\usepackage{courier}  % Font Courier cho monospace

% ── Layout & màu sắc ─────────────────────────────────────────────────────────
\usepackage[top=2.5cm, bottom=2.5cm, left=3cm, right=2.5cm]{geometry}
\usepackage{xcolor}
\definecolor{primary}{RGB}{30,80,160}
\definecolor{codebg}{RGB}{248,248,248}

% ── Tiêu đề section ──────────────────────────────────────────────────────────
\usepackage{titlesec}
\titleformat{\section}{\large\bfseries\color{primary}}{\thesection.}{0.5em}{}[\titlerule]
\titleformat{\subsection}{\normalsize\bfseries\color{primary}}{\thesubsection.}{0.5em}{}
\titleformat{\subsubsection}{\normalsize\bfseries}{\thesubsubsection.}{0.5em}{}

% ── Các gói tiện ích ─────────────────────────────────────────────────────────
\usepackage{graphicx}
\usepackage{hyperref}
\usepackage{enumitem}
\usepackage{array}
\usepackage{booktabs}
\usepackage{listings}
\usepackage{fancyhdr}
\usepackage{amsmath}
\usepackage{float}

\hypersetup{colorlinks=true, linkcolor=primary, urlcolor=primary}

% ── Header / footer ──────────────────────────────────────────────────────────
\pagestyle{fancy}
\fancyhf{}
\rfoot{\thepage}
\renewcommand{\headrulewidth}{0.4pt}

% ── Code listing ─────────────────────────────────────────────────────────────
\lstset{
  backgroundcolor=\color{codebg},
  basicstyle=\ttfamily\small,
  breaklines=true,
  frame=single,
  rulecolor=\color{gray!40},
  numbers=left,
  numberstyle=\tiny\color{gray},
  xleftmargin=1em
}

% =============================================================================
\begin{document}
% =============================================================================

\section{Tổng quan Dự án}

HaNoi Go là một nền tảng web phục vụ khách du lịch nước ngoài đến Hà Nội, tập
trung vào bốn vấn đề cốt lõi:

\begin{enumerate}
  \item \textbf{Khám phá địa điểm:} Thiếu nền tảng tổng hợp thông tin địa điểm
        du lịch Hà Nội có bản đồ tích hợp, phù hợp với khách nước ngoài.
  \item \textbf{Lập lịch trình tối ưu:} Khách du lịch thường mất nhiều thời gian
        di chuyển không cần thiết do không biết cách nhóm địa điểm theo khu vực
        và giờ mở cửa.
  \item \textbf{Kết nối hoạt động:} Khách solo khó tìm người cùng tham gia một
        hoạt động cụ thể trong thời gian ngắn tại Hà Nội.
  \item \textbf{Gợi ý lưu trú và ăn uống:} Khách nước ngoài thiếu công cụ gợi ý
        hotel và nhà hàng phù hợp với lịch trình và khu vực đang ở.
        (Module này được triển khai tuỳ theo tiến độ.)
\end{enumerate}

% =============================================================================
\section{Module 0: User \& Security}
% =============================================================================

Module này xử lý toàn bộ vòng đời tài khoản người dùng --- từ đăng ký, xác
thực email qua OTP, đăng nhập, lấy lại mật khẩu, cập nhật hồ sơ cá nhân, đến
quản trị hệ thống bởi Admin. Đây là nền tảng bảo mật cho tất cả module còn lại,
vì mọi tính năng (lưu lịch trình, tạo activity, chat nhóm) đều yêu cầu xác
thực.

% -------------------------------------------------------------------------
\subsection{Đăng ký tài khoản (Sign Up + OTP Verification)}

\subsubsection{Bài toán}

Khách du lịch nước ngoài cần tạo tài khoản nhanh chóng để lưu lịch trình và
tham gia activity. Form đăng ký phải tối giản, rõ ràng và thân thiện với người
dùng quốc tế. Để tránh tài khoản spam, hệ thống yêu cầu xác thực email bằng
mã OTP trước khi kích hoạt.

\subsubsection{Luồng hoạt động}

\begin{enumerate}
  \item User truy cập trang \texttt{/register}, nhập \textbf{email},
        \textbf{username}, và \textbf{mật khẩu}.
  \item Frontend validate inline: email đúng định dạng, username không chứa
        ký tự đặc biệt, mật khẩu tối thiểu 8 ký tự.
  \item Backend kiểm tra email và username chưa tồn tại trong DB
        (\texttt{UNIQUE} constraint).
  \item Mật khẩu được hash bằng \textbf{bcrypt} (cost factor 12) trước khi
        lưu --- không bao giờ lưu plain text.
  \item Hệ thống tạo bản ghi trong bảng \texttt{users} với
        \texttt{role = USER}, \texttt{status = PENDING}, đồng thời sinh
        \textbf{mã OTP 6 chữ số ngẫu nhiên} có thời hạn 10 phút, lưu vào
        các trường \texttt{otp\_code} và \texttt{otp\_expires\_at}.
  \item Gửi email chứa mã OTP qua \textbf{Nodemailer} (Gmail SMTP).
  \item User được redirect đến trang \texttt{/verify-email?email=...} để
        nhập mã OTP.
  \item \textbf{Xác thực OTP:} Backend kiểm tra mã OTP khớp và chưa hết hạn.
        Nếu hợp lệ, cập nhật \texttt{status = ACTIVE}, xoá \texttt{otp\_code}
        và \texttt{otp\_expires\_at}. Trả về \textbf{Access Token} và
        \textbf{Refresh Token} ngay sau đó --- user không cần đăng nhập lại.
  \item Redirect đến trang \texttt{/profile/edit} để hoàn thiện hồ sơ.
\end{enumerate}

\subsubsection{Tính năng UI}

\begin{itemize}
  \item \textbf{Form đăng ký (\texttt{/register}):} Ba trường Email, Username,
        Password với label rõ ràng. Nút ``Sign Up'' vô hiệu hoá khi đang gửi
        request.
  \item \textbf{Inline validation:} Thông báo lỗi xuất hiện ngay dưới từng
        trường khi blur (rời khỏi ô nhập), không đợi submit.
  \item \textbf{Password strength indicator:} Thanh hiển thị độ mạnh mật
        khẩu (weak / fair / strong) dựa trên độ dài và ký tự đặc biệt.
  \item \textbf{Trang Verify Email (\texttt{/verify-email}):} Giao diện nhập
        OTP gồm 6 ô riêng biệt hỗ trợ tự động focus chuyển ô, paste toàn bộ
        mã, và đếm ngược 60 giây trước khi cho phép ``Resend''. Nếu OTP sai
        hoặc hết hạn, hiển thị thông báo lỗi và user có thể gửi lại mã.
  \item \textbf{Link chuyển trang:} ``Already have an account? Log in'' dẫn
        đến trang Login.
  \item \textbf{Error toasts:} Hiển thị lỗi từ server (email đã tồn tại,
        username đã dùng) dưới dạng toast notification.
\end{itemize}

% -------------------------------------------------------------------------
\subsection{Đăng nhập (Login)}

\subsubsection{Bài toán}

User đã có tài khoản cần đăng nhập để truy cập các tính năng cần xác thực.
Hệ thống dùng \textbf{JWT stateless} --- không lưu session trên server, phù
hợp với kiến trúc NestJS + Next.js tách biệt.

\subsubsection{Luồng hoạt động}

\begin{enumerate}
  \item User nhập \textbf{email hoặc username} và \textbf{mật khẩu} tại trang
        \texttt{/login}. Backend tìm kiếm theo email trước, nếu không có thì
        thử theo username.
  \item Backend xác thực: so sánh mật khẩu với hash bằng
        \texttt{bcrypt.compare()}.
  \item Kiểm tra \texttt{status}: nếu \texttt{BANNED}, trả về lỗi
        \texttt{403} kèm thông báo tài khoản bị khoá. Nếu \texttt{PENDING},
        trả về lỗi yêu cầu xác thực email.
  \item Nếu hợp lệ, sinh \textbf{Access Token} (JWT, TTL 7 ngày, chứa
        \texttt{sub}, \texttt{username}, \texttt{role}, \texttt{tokenVersion})
        và \textbf{Refresh Token} (JWT, TTL 7 ngày, lưu trong
        \texttt{httpOnly cookie}).
  \item Frontend lưu Access Token trong memory (Zustand store). Khi Access
        Token hết hạn, frontend gọi endpoint \texttt{POST /auth/refresh} để
        lấy token mới mà không cần user đăng nhập lại.
  \item Trường \texttt{tokenVersion} trong DB dùng để invalidate token ngay
        lập tức khi user bị ban --- nếu \texttt{tokenVersion} trong JWT không
        khớp với DB, token bị từ chối.
\end{enumerate}

\subsubsection{Tính năng UI}

\begin{itemize}
  \item \textbf{Form đăng nhập:} Hai trường Email/Username và Password.
        Nút ``Log In'' có loading spinner khi đang xử lý.
  \item \textbf{Show / hide password:} Icon con mắt để toggle hiển thị
        mật khẩu.
  \item \textbf{``Forgot password?'':} Link dẫn đến trang Forgot Password.
  \item \textbf{Redirect sau login:} Nếu user bị chặn do chưa đăng nhập
        (ví dụ cố vào \texttt{/trips}), sau khi login thành công sẽ
        redirect về trang đó.
  \item \textbf{Error state:} Thông báo ``Tài khoản hoặc mật khẩu không
        đúng'' khi sai thông tin, không tiết lộ chi tiết (tránh user
        enumeration).
\end{itemize}

% -------------------------------------------------------------------------
\subsection{Đặt lại mật khẩu (Forgot / Reset Password)}

\subsubsection{Bài toán}

User quên mật khẩu cần có cơ chế lấy lại tài khoản an toàn qua email.
Token reset có thời hạn ngắn, chỉ dùng được một lần, và được thiết kế theo
cơ chế \textbf{stateless} --- không cần bảng riêng trong DB.

\subsubsection{Luồng hoạt động}

\begin{enumerate}
  \item \textbf{Bước 1 --- Yêu cầu reset:} User vào trang
        \texttt{/forgot-password}, nhập địa chỉ email.
  \item Backend kiểm tra email có tồn tại trong DB. Nếu không tồn tại,
        vẫn trả về thông báo thành công (tránh lộ thông tin tài khoản).
  \item Hệ thống sinh \textbf{reset token} bằng JWT, sử dụng
        \textbf{secret động} = \texttt{JWT\_SECRET + ":" + passwordHash}
        hiện tại. Token có TTL 15 phút.
  \item Gửi email chứa link dạng
        \texttt{/reset-password?token=\textless token\textgreater}
        qua \textbf{Nodemailer} (Gmail SMTP).
  \item \textbf{Bước 2 --- Đặt mật khẩu mới:} User click link trong
        email, được dẫn đến trang \texttt{/reset-password}.
  \item Backend verify token: giải mã để lấy \texttt{sub} (userId), tìm
        user, tái tạo secret và xác minh chữ ký. Cơ chế này đảm bảo
        tính \textbf{one-time use} vì sau khi mật khẩu thay đổi,
        \texttt{passwordHash} thay đổi → secret cũ không còn hợp lệ.
  \item Mật khẩu mới được hash bcrypt và cập nhật vào bảng \texttt{users}.
\end{enumerate}

\subsubsection{Tính năng UI}

\begin{itemize}
  \item \textbf{Trang Forgot Password (\texttt{/forgot-password}):}
        Ô nhập email, nút ``Send Reset Link''. Sau khi gửi, hiển thị
        thông báo ``Check your inbox'' dù email có tồn tại hay không.
  \item \textbf{Trang Reset Password (\texttt{/reset-password}):}
        Hai ô nhập ``New Password'' và ``Confirm Password'' với kiểm tra
        khớp nhau. Nếu token hết hạn hoặc không hợp lệ, hiển thị thông
        báo lỗi và nút ``Request a new link''.
  \item \textbf{Thành công:} Sau khi đặt lại mật khẩu, redirect về trang
        Login với toast ``Mật khẩu đã được cập nhật thành công.''.
\end{itemize}

% -------------------------------------------------------------------------
\subsection{Hồ sơ người dùng (Profile)}

\subsubsection{Bài toán}

Sau khi xác thực, user cần một trang hồ sơ cá nhân để xem lại thống kê
hoạt động (số trip, địa điểm yêu thích, activities đã tham gia) và chỉnh
sửa thông tin. Trang này cũng phục vụ cho khách du lịch nước ngoài muốn
khai báo quốc tịch và ngôn ngữ để cộng đồng dễ kết nối.

\subsubsection{Mô hình dữ liệu người dùng (User)}

\begin{lstlisting}[language=SQL]
CREATE TABLE users (
  id            uuid PRIMARY KEY DEFAULT gen_random_uuid(),
  email         text UNIQUE NOT NULL,
  password_hash text NOT NULL,
  username      text UNIQUE NOT NULL,
  full_name     text,
  bio           text,
  avatar_url    text,
  nationality   text,
  languages     text[],           -- ['English', 'French', ...]
  role          text DEFAULT 'USER',    -- USER | ADMIN
  status        text DEFAULT 'PENDING', -- PENDING | ACTIVE | BANNED
  token_version int  DEFAULT 1,        -- invalidate JWT khi ban
  otp_code      text,                  -- 6-digit OTP
  otp_expires_at timestamptz,
  created_at    timestamptz DEFAULT now()
);
\end{lstlisting}

\subsubsection{Tính năng Profile UI (\texttt{/profile})}

\begin{itemize}
  \item \textbf{Hero header:} Ảnh bìa panorama Hà Nội, avatar tròn nổi
        phía trên, tên đầy đủ và username.
  \item \textbf{Thông tin cá nhân:} Bio, quốc tịch (\texttt{nationality}),
        danh sách ngôn ngữ (\texttt{languages[]}).
  \item \textbf{Quick Stats Grid:} 4 ô thống kê --- Trips (số lịch trình
        đã tạo), Places (địa điểm đã lưu trong Basket), Activities (số
        activity đã tham gia), Languages.
  \item \textbf{3 Tab nội dung:}
        \begin{itemize}
          \item \textbf{Trips:} Danh sách lịch trình đã tạo, mỗi card
                hiển thị tiêu đề, số ngày, ngày tạo và link ``Open Planner''.
          \item \textbf{Saved:} Danh sách địa điểm đã lưu vào Travel Basket
                (Zustand store), có nút xoá từng địa điểm.
          \item \textbf{Activities:} Danh sách activity đã tham gia, hiển
                thị tiêu đề, ngày giờ, địa chỉ và nút ``View Chat''.
        \end{itemize}
  \item \textbf{Trang Edit Profile (\texttt{/profile/edit}):} Form chỉnh
        sửa \texttt{fullName}, \texttt{bio}, \texttt{nationality},
        \texttt{languages}, upload avatar (qua Media module).
\end{itemize}

% -------------------------------------------------------------------------
\subsection{Admin Site --- Quản trị hệ thống}

\subsubsection{Bài toán}

Hệ thống cần giao diện quản trị toàn diện cho phép Admin theo dõi số liệu
tổng quan, quản lý tài khoản, quản lý địa điểm, và xử lý báo cáo vi phạm.
Admin Site chỉ truy cập được bởi tài khoản có \texttt{role = ADMIN} ---
mọi endpoint đều được bảo vệ bởi \textbf{RolesGuard} (NestJS) kết hợp
kiểm tra middleware trên Next.js.

\subsubsection{Phân quyền (RBAC)}

Hệ thống định nghĩa hai role và ba trạng thái tài khoản:

\begin{center}
\begin{tabular}{lll}
  \toprule
  \textbf{Role} & \textbf{Quyền} & \textbf{Ghi chú} \\
  \midrule
  \texttt{USER}  & Dùng toàn bộ tính năng app & Mặc định khi đăng ký \\
  \texttt{ADMIN} & USER + toàn bộ Admin Site   & Gán thủ công trong DB \\
  \bottomrule
\end{tabular}
\end{center}

\begin{center}
\begin{tabular}{ll}
  \toprule
  \textbf{Status} & \textbf{Ý nghĩa} \\
  \midrule
  \texttt{PENDING} & Chưa xác thực email, chưa đăng nhập được \\
  \texttt{ACTIVE}  & Tài khoản hoạt động bình thường \\
  \texttt{BANNED}  & Bị khoá, trả về \texttt{403} khi cố đăng nhập \\
  \bottomrule
\end{tabular}
\end{center}

\subsubsection{Dashboard (\texttt{/admin/dashboard})}

Trang tổng quan hiển thị các số liệu realtime:

\begin{itemize}
  \item \textbf{Stats Cards:} Tổng số user (role=USER), user bị ban, tổng
        địa điểm, tổng lịch trình, báo cáo đang chờ xử lý.
  \item \textbf{User Growth Chart (30 ngày):} Biểu đồ đường hiển thị số
        tài khoản mới đăng ký theo ngày, truy vấn trực tiếp từ PostgreSQL
        (group by DATE).
  \item \textbf{Violation Breakdown:} Biểu đồ tròn phân loại vi phạm theo
        \texttt{ViolationType} (SPAM, HATE\_SPEECH, SCAM, FAKE\_INFORMATION,
        INAPPROPRIATE\_CONTENT, OTHER).
  \item \textbf{Report Status Summary:} Số báo cáo theo từng trạng thái
        (PENDING / RESOLVED / DISMISSED).
  \item \textbf{Popular Places:} Top 5 địa điểm được ghé thăm nhiều nhất
        (đếm qua bảng \texttt{trip\_stops}).
\end{itemize}

\subsubsection{Quản lý người dùng (\texttt{/admin/users})}

\begin{itemize}
  \item \textbf{Danh sách user:} Bảng dữ liệu với các cột: Avatar, Username,
        Email, Full Name, Role, Status, Ngày tạo, Trip Count. Hỗ trợ
        \textbf{phân trang} (10 user/trang), \textbf{sắp xếp} và
        \textbf{lọc theo status}.
  \item \textbf{Tìm kiếm:} Thanh search lọc theo email, username hoặc
        fullName (tìm kiếm không phân biệt hoa thường, \texttt{ILIKE}).
  \item \textbf{Khoá tài khoản (Ban):} Chuyển \texttt{status} sang
        \texttt{BANNED}, tăng \texttt{tokenVersion} (invalidate JWT ngay lập
        tức), tạo bản ghi \texttt{Report} và \texttt{AdminLog} trong
        transaction.
  \item \textbf{Mở khoá tài khoản (Unban):} Chuyển \texttt{status} về
        \texttt{ACTIVE}, tạo \texttt{AdminLog}.
  \item \textbf{Tạo user thủ công:} Admin có thể tạo tài khoản trực tiếp
        (bao gồm cả tài khoản ADMIN mới), bỏ qua bước OTP.
  \item \textbf{Xoá user:} Xoá hoàn toàn, cascade xoá reports và adminLogs
        liên quan trong transaction.
  \item \textbf{Xem chi tiết user:} Trang profile đọc-only hiển thị thông
        tin tài khoản, số trips, số báo cáo đã tạo và bị báo cáo.
\end{itemize}

\subsubsection{Quản lý địa điểm (\texttt{/admin/places})}

\begin{itemize}
  \item \textbf{Danh sách địa điểm:} Bảng hiển thị tên, category, district,
        số lần xuất hiện trong \texttt{trip\_stops}. Hỗ trợ tìm kiếm và lọc
        theo category.
  \item \textbf{Thêm địa điểm mới:} Form nhập đầy đủ thông tin (tên, toạ độ,
        category, district, giờ mở cửa, ảnh). Backend dùng
        \texttt{ST\_GeomFromEWKT} để lưu geometry PostGIS.
  \item \textbf{Chỉnh sửa địa điểm:} Cập nhật mọi trường, bao gồm cả toạ độ
        (cập nhật đồng thời cột \texttt{lat}, \texttt{lng} và
        \texttt{location} geometry).
  \item \textbf{Xoá địa điểm:} Xoá địa điểm, cascade xoá gallery ảnh.
\end{itemize}

\subsubsection{Quản lý báo cáo vi phạm (\texttt{/admin/reports})}

\begin{itemize}
  \item \textbf{Danh sách báo cáo:} Bảng hiển thị người báo cáo, đối tượng
        bị báo cáo, loại vi phạm, trạng thái, ngày tạo. Lọc theo
        \texttt{ReportStatus} (PENDING / RESOLVED / DISMISSED) và tìm kiếm
        theo username.
  \item \textbf{Chi tiết báo cáo:} Xem thông tin đầy đủ, nếu vi phạm gắn với
        Activity (\texttt{entityType = ACTIVITY}) thì hiển thị thông tin
        activity đó kèm host.
  \item \textbf{Xử lý báo cáo:} Admin ghi chú (\texttt{adminNotes}), đánh dấu
        RESOLVED (tuỳ chọn ẩn activity bị báo cáo bằng cách chuyển sang
        \texttt{CANCELLED}) hoặc DISMISSED.
  \item \textbf{AdminLog:} Mọi hành động Admin (BAN\_USER, UNBAN\_USER,
        RESOLVE\_REPORT, DISMISS\_REPORT) đều được ghi vào bảng
        \texttt{admin\_logs} để audit trail.
\end{itemize}

\subsubsection{Bảo vệ Route}

\begin{lstlisting}[language={}]
-- Backend: RolesGuard (NestJS)
@Roles('ADMIN')
@UseGuards(JwtAuthGuard, RolesGuard)
@Get('/admin/users')
listUsers() { ... }

-- Frontend: Middleware Next.js
-- Redirect ve /login neu role !== 'ADMIN'
export function middleware(request: NextRequest) {
  const token = request.cookies.get('access_token');
  const payload = decodeJwt(token);
  if (payload.role !== 'ADMIN') {
    return NextResponse.redirect('/login');
  }
}
\end{lstlisting}

% =============================================================================
\section{Các Module Chính}
% =============================================================================

\subsection{Khám phá Địa điểm (Discovery)}

\subsubsection{Bài toán}

Khách du lịch nước ngoài cần một nơi tra cứu thông tin địa điểm du lịch tại
Hà Nội một cách trực quan, có bản đồ tích hợp và khả năng lọc theo khu vực
hoặc loại hình. Module này đóng vai trò \textbf{foundation data} cho toàn bộ
hệ thống --- bộ dữ liệu 65 địa điểm du lịch đã được xác thực, sẵn sàng hỗ trợ
tìm đường thời gian thực.

\subsubsection{Tính năng}

\begin{itemize}
  \item \textbf{Danh sách địa điểm (\texttt{/discovery}):} Hiển thị tên, ảnh,
        loại hình (category), khu vực (district), giờ mở cửa, thời gian tham
        quan ước tính và tags.
  \item \textbf{Filter:} Lọc theo khu vực (Hoàn Kiếm, Ba Đình, Tây Hồ, Đống
        Đa, Hai Bà Trưng, Cầu Giấy) và loại hình (Heritage \& History,
        Arts \& Culture, Sightseeing, Eat \& Shop, Nature \& Outdoors,
        Spiritual).
  \item \textbf{Map view:} Hiển thị toàn bộ địa điểm dưới dạng pin trên bản
        đồ Leaflet (CartoDB Voyager tiles). Click vào pin để xem thông tin
        chi tiết.
  \item \textbf{Trang chi tiết địa điểm (\texttt{/places/:id}):} Thông tin
        đầy đủ kèm vị trí trên bản đồ, gallery ảnh (\texttt{PlaceGallery}),
        và địa điểm lân cận trong bán kính 1\,km (tái dụng
        \texttt{ST\_DWithin} PostGIS).
  \item \textbf{Travel Basket:} User có thể lưu địa điểm yêu thích vào
        \texttt{TravelBasket} (Zustand store, persist localStorage). Basket
        là đầu vào cho Trip Planner.
\end{itemize}

\subsubsection{Triển khai}

Dữ liệu địa điểm được lưu trong PostgreSQL với PostGIS. Giờ mở cửa được
tách thành các cột riêng biệt giúp thuật toán đọc nhanh và truy vấn hiệu
quả:

\begin{lstlisting}[language=SQL]
CREATE TABLE places (
  id                uuid PRIMARY KEY DEFAULT gen_random_uuid(),
  name              text NOT NULL UNIQUE,
  description_en    text,
  category          text NOT NULL,
  district          text NOT NULL,
  address           text,
  lat               float NOT NULL,
  lng               float NOT NULL,
  location          geometry(Point, 4326),
  visit_duration_min int,
  image_url         text,
  tags              text[],
  -- Lịch hoạt động (parsed từ opening hours)
  always_open       boolean DEFAULT false,
  open_days         integer[],  -- [0..6] = [CN,T2..T7]
  open_time_start   time,       -- giờ mở cửa (VD: 08:00)
  open_time_end     time,       -- giờ đóng cửa (VD: 17:00)
  has_break         boolean DEFAULT false,
  break_start       time,       -- nghỉ trưa bắt đầu
  break_end         time,       -- nghỉ trưa kết thúc
  created_at        timestamptz DEFAULT now()
);

-- Bảng gallery ảnh (1 địa điểm - nhiều ảnh)
CREATE TABLE place_gallery (
  id        uuid PRIMARY KEY DEFAULT gen_random_uuid(),
  place_id  uuid REFERENCES places(id) ON DELETE CASCADE,
  url       text NOT NULL,
  created_at timestamptz DEFAULT now()
);

-- GiST index để tăng tốc ST_DWithin
CREATE INDEX ON places USING GIST (location);
\end{lstlisting}

% -------------------------------------------------------------------------
\subsection{Trip Planner}

\subsubsection{Bài toán}

Người dùng chọn danh sách địa điểm du lịch muốn đến tại Hà Nội cùng số ngày
và điểm xuất phát. Hệ thống tự động sắp xếp lịch trình theo ngày sao cho tổng
quãng đường di chuyển là nhỏ nhất, đồng thời thoả mãn ràng buộc về giờ mở cửa
và thời gian tham quan. Lịch trình được hiển thị trực tiếp trên bản đồ Leaflet.

\medskip
\textbf{Tình huống minh hoạ:} Một du khách nước ngoài đang đứng tại khu vực
Nhà hát Lớn muốn thăm 5 địa điểm trong 1 ngày (08:00--18:00): Hồ Hoàn Kiếm,
Cha Cá Lã Vọng, Đền Bạch Mã, Ca Trù Thăng Long, và Văn Miếu. Mục tiêu là hệ
thống tự phân bổ thứ tự hợp lý nhất --- không để khách đi vòng ziczac giữa hai
đầu phố, không xếp lịch vào địa điểm đang đóng cửa.

\subsubsection{Thuật toán tổng quan}

Bài toán thuộc dạng biến thể của \textbf{Travelling Salesman Problem with Time
Windows (TSPTW)} mở rộng sang đa ngày. Hệ thống giải quyết theo pipeline
\textbf{Cluster-first, Route-second} gồm 5 giai đoạn nối tiếp nhau.

% ── Giai đoạn 1 ─────────────────────────────────────────────────────────────
\paragraph{Giai đoạn 1 --- Lọc và Phân cụm theo Địa lý (K-Means)}

\textit{Kỹ thuật:} Trước tiên, hệ thống tra cứu địa điểm theo \texttt{placeIds}
(UUID) từ cơ sở dữ liệu, lọc bỏ các địa điểm đóng cửa hoàn toàn trong toàn bộ
hành trình. Các địa điểm còn lại được phân vào $D$ cụm (một cụm = một ngày)
bằng thuật toán K-Means với khoảng cách Haversine:
\[
  d = 2R \arcsin\!\sqrt{
    \sin^2\!\tfrac{\Delta\phi}{2}
    + \cos\phi_1\cos\phi_2\sin^2\!\tfrac{\Delta\lambda}{2}
  }
\]
Tâm cụm ban đầu được chọn theo \textbf{K-Means++} (xác suất tỉ lệ khoảng cách)
kết hợp \textbf{Seeded PRNG} để đảm bảo cùng một bộ địa điểm luôn cho ra kết
quả phân cụm như nhau. Sau K-Means, nếu một địa điểm rơi vào ngày mà nó đóng
cửa, bước \texttt{postClusterOpenDaySwap} sẽ hoán đổi nó sang cụm phù hợp hơn
--- mà không làm thay đổi chỉ số ngày của các cụm khác.

\medskip
\textit{Tình huống thực tế:} Với 5 địa điểm trên, thuật toán nhận thấy
Văn Miếu nằm ở phía Đống Đa, cách xa 4 địa điểm còn lại ở Hoàn Kiếm. K-Means
tự động tách Văn Miếu sang Ngày 2, còn Ngày 1 gom 4 địa điểm Hoàn Kiếm lại với
nhau --- du khách không phải đi ngược từ Hoàn Kiếm sang Đống Đa rồi quay lại.

% ── Giai đoạn 2 ─────────────────────────────────────────────────────────────
\paragraph{Giai đoạn 2 --- Tính toán Cửa sổ Thời gian (\texttt{calculateVisitWindow})}

\textit{Kỹ thuật:} Đây là hàm trung tâm, được gọi nhất quán ở mọi giai đoạn
(GNN, cascade, conflict resolution). Với mỗi địa điểm $p$ và thời điểm dự kiến
đến $T_{arrival}$, hàm tính:
\begin{enumerate}
  \item Điều chỉnh theo giờ mở cửa: $T_{visit} = \max(T_{arrival},\, O_{start})$.
  \item \textbf{Multi-pass blocked windows:} Thu thập tất cả khoảng ``cấm''
        (giờ nghỉ trưa chung $[L_{start}, L_{end}]$ và giờ nghỉ riêng của địa
        điểm $[B_{start}, B_{end}]$ nếu có), sắp xếp theo thứ tự. Nếu khoảng
        tham quan $[T_{visit},\, T_{visit}+v]$ chồng lên bất kỳ khoảng cấm nào,
        đẩy $T_{visit}$ ra sau khoảng đó và kiểm tra lại từ đầu.
  \item Kiểm tra hợp lệ: $T_{visit}+v \le \min(O_{end},\, T_{end})$. Nếu vi
        phạm, trả về \texttt{null}.
\end{enumerate}

\medskip
\textit{Tình huống thực tế:} Du khách đến Cha Cá Lã Vọng lúc 10:43, tham quan
60 phút. Khoảng $[10{:}43, 11{:}43]$ chồng vào giờ nghỉ trưa $[11{:}00, 13{:}00]$.
Hàm phát hiện xung đột, đẩy $T_{visit}$ về 13:00 và kiểm tra lại --- không còn
chồng chéo nào, trả về \textit{bắt đầu 13:00, kết thúc 14:00}. Cách xử lý
multi-pass này đảm bảo dù địa điểm có thêm giờ nghỉ riêng buổi chiều, logic vẫn
hoạt động đúng mà không cần thêm code đặc biệt.

% ── Giai đoạn 3 ─────────────────────────────────────────────────────────────
\paragraph{Giai đoạn 3 --- Định tuyến Tham lam có Ràng buộc Thời gian (Time-Window GNN)}

\textit{Kỹ thuật:} Trong mỗi cụm, hệ thống gọi \textbf{Goong Distance Matrix
API} một lần duy nhất để lấy ma trận $N \times N$ thời gian di chuyển thực tế
bằng xe máy giữa tất cả các cặp địa điểm. Ma trận này được lưu lại để tái sử
dụng ở các giai đoạn sau. Thuật toán GNN sau đó chọn điểm tiếp theo theo hàm
chi phí:
\begin{equation}
  Cost(j) = 2 \times \Delta T(i \to j) + T_{wait}(j) \times 60
\end{equation}
Hệ số $\times 2$ phạt nặng di chuyển xa hơn thời gian chờ, đồng thời phá thế
cân bằng khi nhiều điểm cùng bị đẩy ra sau giờ nghỉ trưa (tất cả đều bắt đầu
lúc 13:00 nhưng cách nhau về vị trí địa lý).

\medskip
\textit{Tình huống thực tế:} Lúc 10:43 sau khi rời Cha Cá Lã Vọng, thuật toán
so sánh hai ứng viên: Đền Bạch Mã (200 m, đến 10:44) và Ca Trù Thăng Long
(1,2 km, đến 10:55). Cả hai đều bị đẩy sang 13:00 vì giờ nghỉ trưa, nên
$T_{wait}$ bằng nhau. Khi đó $\Delta T \times 2$ phân biệt: Đền Bạch Mã có chi
phí thấp hơn → chọn Đền Bạch Mã đi trước. Kết quả là Ca Trù được xếp ngay
sau đó lúc 14:10, không cần đi ngược lại.

% ── Giai đoạn 4 ─────────────────────────────────────────────────────────────
\paragraph{Giai đoạn 4 --- Điều chỉnh GPS Thực tế (\texttt{cascadeRouteTimes})}

\textit{Kỹ thuật:} GNN chạy trên giả định xuất phát từ 08:00 không có thời gian
di chuyển ban đầu. Sau đó, nếu người dùng cung cấp tọa độ GPS, hàm
\texttt{cascadeRouteTimes} áp dụng offset thực tế:
\begin{itemize}
  \item \textbf{Điểm đầu tiên:} Gọi Goong API để tính thời gian từ GPS người
        dùng đến điểm đầu; toàn bộ lịch trình dịch chuyển theo.
  \item \textbf{Các điểm tiếp theo:} Tái sử dụng \texttt{travelFromPrevSec} đã
        lưu từ Goong matrix ở Giai đoạn 3 --- không tính lại bằng Haversine.
  \item \textbf{Khi có drop:} Nếu một điểm vi phạm ràng buộc sau khi cascade,
        nó bị loại và thời gian di chuyển của điểm kế tiếp được tính lại qua
        Goong API (vì predecessor đã thay đổi).
\end{itemize}

\medskip
\textit{Tình huống thực tế:} GNN ban đầu xếp Hồ Hoàn Kiếm lúc 08:00. Nhưng
du khách đang đứng ở Nhà hát Lớn --- Goong API cho biết đi xe máy mất 18 phút.
Cascade dịch toàn bộ lịch: Hồ Hoàn Kiếm dời thành 08:28. Điểm kế tiếp (Cha Cá
Lã Vọng) dùng lại thời gian di chuyển đã biết từ ma trận (5 phút) → đến 09:43.
Không có địa điểm nào vi phạm giờ đóng cửa sau khi dịch, nên không có drop.

% ── Giai đoạn 5 ─────────────────────────────────────────────────────────────
\paragraph{Giai đoạn 5 --- Giải quyết Xung đột (Gap Insertion)}

\textit{Kỹ thuật:} Nếu sau GNN + cascade còn địa điểm chưa được xếp lịch (do
cụm ngày đó quá chật), hệ thống thử chèn từng địa điểm đó vào mọi \textit{khe
trống} (gap) của tất cả các ngày. Với mỗi vị trí chèn $k$ trong ngày $d$:
\begin{align}
  T_{arrival}(p) &= T_{depart}(k-1) + \Delta T(k-1, p) \\
  T_{depart}(p)  &= \texttt{calculateVisitWindow}(p,\; T_{arrival}(p))\texttt{.departMin}
\end{align}
Điều kiện chèn thành công: (1) \texttt{calculateVisitWindow} trả về khác
\texttt{null}, và (2) $T_{depart}(p) + \Delta T(p, k) \le T_{arrival}(k)$
(không làm trễ điểm đứng sau). Nếu không tìm được vị trí nào trong toàn bộ
hành trình, địa điểm đi vào \texttt{unscheduled[]} và Frontend hiển thị thông
báo lý do.

\subsubsection{Tính toán Khoảng cách và API}

Hệ thống sử dụng hai lớp khoảng cách với mục đích khác nhau:

\begin{table}[H]
\centering
\renewcommand{\arraystretch}{1.3}
\begin{tabular}{@{} p{0.18\textwidth} p{0.22\textwidth} p{0.50\textwidth} @{}}
\toprule
\textbf{Phương pháp} & \textbf{Dùng ở đâu} & \textbf{Lý do} \\
\midrule
Haversine & K-Means clustering & Chỉ cần khoảng cách tương đối để gom cụm địa lý, không cần chính xác theo đường phố. \\
\addlinespace
Goong Distance Matrix API (\texttt{vehicle=bike}) & GNN routing, GPS cascade & Cần thời gian di chuyển thực tế theo đường phố Hà Nội để lịch không sai lệch. \\
\addlinespace
Haversine (fallback) & Khi Goong API thất bại & Đảm bảo hệ thống không crash khi API quá tải. \\
\bottomrule
\end{tabular}
\caption{Phân tầng khoảng cách trong Trip Planner}
\end{table}

\textbf{Xử lý giới hạn tốc độ API:} Khi nhận mã lỗi \texttt{OVER\_RATE\_LIMIT},
hệ thống áp dụng \textbf{Exponential Backoff} --- thử lại 3 lần với độ trễ
tăng dần 1s $\to$ 2s $\to$ 4s trước khi chuyển sang Haversine fallback.

\subsubsection{Cơ sở dữ liệu Trip}

\begin{lstlisting}[language=SQL]
CREATE TABLE trips (
  id            uuid PRIMARY KEY DEFAULT gen_random_uuid(),
  user_id       uuid REFERENCES users(id) ON DELETE CASCADE,
  title         text,
  num_days      int NOT NULL,
  start_place_id uuid REFERENCES places(id),
  cloned_from_id uuid REFERENCES trips(id) ON DELETE SET NULL,
  created_at    timestamptz DEFAULT now()
);

CREATE TABLE trip_days (
  id         uuid PRIMARY KEY DEFAULT gen_random_uuid(),
  trip_id    uuid REFERENCES trips(id) ON DELETE CASCADE,
  day_number int NOT NULL,
  district   text  -- khu vực chủ đạo của ngày
);

CREATE TABLE trip_stops (
  id                  uuid PRIMARY KEY DEFAULT gen_random_uuid(),
  trip_day_id         uuid REFERENCES trip_days(id) ON DELETE CASCADE,
  place_id            uuid REFERENCES places(id),
  stop_order          int NOT NULL,
  arrive_at           time,
  depart_at           time,
  distance_from_prev_m int,
  duration_from_prev_s  int,
  is_skipped          boolean DEFAULT false
);
\end{lstlisting}

Trường \texttt{cloned\_from\_id} cho phép truy vết nguồn gốc khi lịch trình
được sao chép từ Shared Trips.

\subsubsection{Ví dụ Đầu vào / Đầu ra}

\textit{Bối cảnh:} Du khách đang ở Nhà hát Lớn (21.0285, 105.8412), muốn thăm
4 địa điểm tại Hoàn Kiếm trong ngày 13/05/2026, từ 08:00 đến 18:00.

\begin{lstlisting}[language={}]
-- INPUT --
placeIds: ["uuid-hoan-kiem", "uuid-cha-ca", "uuid-bach-ma", "uuid-ca-tru"]
numDays: 1, startTime: 480 (08:00), endTime: 1080 (18:00)
travelDate: "2026-05-13", startLat: 21.0285, startLng: 105.8412

-- OUTPUT --
Ngay 1:
  [18 phut tu GPS] -> Ho Hoan Kiem    08:28 - 09:28
  [ 5 phut di chuyen]  Cha Ca La Vong 09:43 - 10:43
  [Nghi trua 11:00 - 13:00]
  [ 1 phut di chuyen]  Den Bach Ma   13:01 - 14:01
  [ 8 phut di chuyen]  Ca Tru TL     14:10 - 15:10

unscheduled: []   <-- tat ca deu duoc xep lich
\end{lstlisting}

\subsubsection{Đánh giá và Tinh chỉnh Thuật toán}

\begin{table}[htbp]
\centering
\caption{Bảng so sánh tối ưu hoá thuật toán Trip Planner}
\label{tab:algorithm_optimization}
\renewcommand{\arraystretch}{1.4}
\begin{tabular}{@{} p{0.20\textwidth} p{0.36\textwidth} p{0.36\textwidth} @{}}
\toprule
\textbf{Thành phần} & \textbf{Phiên bản cơ sở (Baseline)} & \textbf{Phiên bản tinh chỉnh (Refactored)} \\
\midrule
\textbf{Tra cứu địa điểm} & Tra cứu theo tên (\texttt{placeNames}). Lỗi nếu tên trùng hoặc đổi tên. & Tra cứu theo UUID (\texttt{placeIds}), hỗ trợ tên như fallback tương thích ngược. \\
\addlinespace
\textbf{Khởi tạo K-Means} & Chọn tâm cụm theo chỉ số mảng tuần tự. Gom cụm kém nếu dữ liệu không có thứ tự địa lý. & Áp dụng \textbf{K-Means++} kết hợp \textbf{Seeded PRNG} tạo tính xác định (Deterministic). \\
\addlinespace
\textbf{Bảo toàn chỉ số cụm} & Lọc bỏ cụm rỗng sau clustering, gây lệch ánh xạ ngày-cụm. & Giữ nguyên cụm rỗng; vòng lặp scheduling bỏ qua ngày trống an toàn. \\
\addlinespace
\textbf{Cửa sổ thời gian} & Xử lý giờ nghỉ rời rạc tại nhiều chỗ, dễ bỏ sót giờ nghỉ riêng của địa điểm. & Hàm \textbf{\texttt{calculateVisitWindow}} tập trung, multi-pass, kiểm tra \texttt{openTimeEnd} nhất quán. \\
\addlinespace
\textbf{Cascade GPS} & Bỏ Goong matrix, tính lại travel times bằng Haversine --- kết quả sai (ví dụ: 1 phút cho 2 điểm cách nhau 500m theo đường phố). & Tái sử dụng \texttt{travelFromPrevSec} từ Goong matrix. Chỉ gọi API lại khi stop bị drop. \\
\addlinespace
\textbf{Giải quyết xung đột} & Chỉ nối điểm rớt lịch vào \textit{cuối ngày}. Gán cứng 10 phút di chuyển. & \textbf{Gap Insertion} quét mọi khoảng trống trong ngày, dùng thời gian thực tế. \\
\addlinespace
\textbf{Dự phòng Rate Limit} & Thử lại 1 lần, chờ 3 giây cố định. & \textbf{Exponential Backoff} 3 lần (1s $\to$ 2s $\to$ 4s). \\
\addlinespace
\textbf{Chèn kẽ stop} & Gọi \texttt{calculateVisitWindow} thiếu thời gian di chuyển và đệm đỗ xe (mặc định 0), gây sai lệch khi chèn stop vào khoảng trống. & Truyền đầy đủ thời gian di chuyển thực tế từ stop liền trước và \texttt{PARKING\_BUFFER\_MIN} vào hàm \texttt{calculateVisitWindow}. \\
\addlinespace
\textbf{Định tuyến ngày đóng} & Chuyển địa điểm đóng sang ngày khác mà không kiểm tra sức chứa tối đa của ngày đích. & Kiểm tra dung lượng ngày đích có nhỏ hơn \texttt{MAX\_PLACES\_PER\_DAY} (5 địa điểm) trước khi tiến hành hoán đổi cụm. \\
\bottomrule
\end{tabular}
\end{table}


% -------------------------------------------------------------------------
\subsection{Activity \& Group Chat}

\subsubsection{Bài toán}

Khách du lịch solo muốn tìm người cùng tham gia một hoạt động cụ thể tại Hà
Nội trong thời gian ngắn --- nhưng các nền tảng hiện tại (Facebook, Meetup)
không gắn hoạt động trực tiếp lên bản đồ và không hỗ trợ group chat tức thì
sau khi ghép nhóm.

\subsubsection{Luồng hoạt động}

\begin{enumerate}
  \item \textbf{Tạo activity:} User điền tên hoạt động, category (Arts \&
        Culture, Heritage \& History, Sightseeing,...), thời gian, mô tả, số
        người tối đa, ảnh bìa. Địa điểm họp mặt được tìm kiếm thông qua cơ
        chế tìm kiếm lai ghép (\textit{hybrid search}): hệ thống ưu tiên đối
        khớp và gợi ý các địa danh du lịch trọng điểm sẵn có trong cơ sở dữ
        liệu (65 địa danh) kèm huy hiệu nổi bật (``Highlight''), đồng thời
        kết hợp \textbf{Goong Autocomplete API} để tìm kiếm các địa điểm
        ngoài danh mục (quán cà phê, khách sạn, v.v.). Khi người dùng chọn
        gợi ý hoặc kéo thả ghim trên bản đồ thu nhỏ (\textbf{LocationPickerMap}),
        hệ thống tự động liên kết với \texttt{place\_id} tương ứng trong cơ
        sở dữ liệu (nếu có), đồng thời kích hoạt bộ chuẩn hoá và dịch địa
        chỉ sang tiếng Anh không dấu (\texttt{translateAddressToEnglish}).
  \item \textbf{Hiển thị trên map:} Activity xuất hiện dưới dạng marker thông
        minh trên bản đồ Leaflet. Trạng thái và màu sắc marker tự động thay
        đổi theo thời gian họp mặt (\texttt{scheduled\_at}):
        \begin{itemize}
          \item \textbf{UPCOMING} (Chưa diễn ra): Marker màu xanh dương
                (\texttt{\#1e50a0}) tiêu chuẩn.
          \item \textbf{STARTING} (Sắp bắt đầu trong 2 giờ): Marker màu vàng
                hổ phách (\texttt{\#f59e0b}) kèm badge hoạt họa nhấp nháy
                ``SOON''.
          \item \textbf{ONGOING} (Đang diễn ra trong vòng 3 giờ từ lúc bắt
                đầu): Marker màu xanh lá cây (\texttt{\#16a34a}) kèm hiệu ứng
                quét sóng tròn (\textit{pulse animation}) và badge ``LIVE''.
          \item \textbf{ENDED} (Đã kết thúc sau 3 giờ): Marker màu xám
                (\texttt{\#9ca3af}) mờ với độ trong suốt 50\%.
        \end{itemize}
  \item \textbf{Duyệt thành viên:} Chủ activity nhận thông báo
        (\textbf{Notification} loại \texttt{ACTIVITY\_REQUEST}), xem profile
        người xin tham gia và duyệt (\texttt{APPROVED}) hoặc từ chối
        (\texttt{REJECTED}). Người xin tham gia nhận thông báo kết quả
        (\texttt{ACTIVITY\_APPROVED} / \texttt{ACTIVITY\_REJECTED}).
  \item \textbf{Group chat:} Sau khi được duyệt, thành viên tự động vào group
        chat của activity. Chat đồng bộ realtime qua WebSocket và Redis
        Pub/Sub. Tin nhắn có hai loại: \texttt{TEXT} (người dùng gửi) và
        \texttt{SYSTEM} (thông báo hệ thống như ``User X đã tham gia nhóm'').
  \item \textbf{Message Reactions:} Thành viên có thể react emoji vào từng
        tin nhắn. Mỗi cặp (message, user, emoji) là duy nhất --- tránh react
        trùng.
  \item \textbf{Báo cáo vi phạm:} User có thể báo cáo activity vi phạm.
        Report được gửi đến Admin Site với \texttt{entityType = ACTIVITY}.
\end{enumerate}

\subsubsection{Triển khai}

\begin{lstlisting}[language=SQL]
CREATE TABLE activities (
  id           uuid PRIMARY KEY DEFAULT gen_random_uuid(),
  host_id      uuid REFERENCES users(id) ON DELETE CASCADE,
  place_id     uuid REFERENCES places(id),  -- null neu tu nhap
  trip_id      uuid REFERENCES trips(id) ON DELETE CASCADE, -- neu la Shared Trip
  title        text NOT NULL,
  description  text,
  address      text,
  lat          float NOT NULL,
  lng          float NOT NULL,
  location     geometry(Point, 4326),
  scheduled_at timestamptz NOT NULL,
  max_members  int DEFAULT 10,
  status       text DEFAULT 'OPEN',  -- OPEN | CLOSED | CANCELLED
  category     text DEFAULT 'Arts & Culture',
  image_url    text,
  saves_count  int DEFAULT 0,  -- dem luot Save Trip
  created_at   timestamptz DEFAULT now()
);

CREATE TABLE activity_members (
  activity_id  uuid REFERENCES activities(id) ON DELETE CASCADE,
  user_id      uuid REFERENCES users(id) ON DELETE CASCADE,
  status       text CHECK (status IN ('PENDING','APPROVED','REJECTED')),
  joined_at    timestamptz DEFAULT now(),
  PRIMARY KEY (activity_id, user_id)
);

CREATE TABLE messages (
  id          uuid PRIMARY KEY DEFAULT gen_random_uuid(),
  activity_id uuid REFERENCES activities(id) ON DELETE CASCADE,
  user_id     uuid REFERENCES users(id) ON DELETE SET NULL,
  content     text NOT NULL,
  type        text DEFAULT 'TEXT',  -- TEXT | SYSTEM
  created_at  timestamptz DEFAULT now()
);

CREATE TABLE message_reactions (
  id          uuid PRIMARY KEY DEFAULT gen_random_uuid(),
  message_id  uuid REFERENCES messages(id) ON DELETE CASCADE,
  user_id     uuid REFERENCES users(id) ON DELETE CASCADE,
  emoji       text NOT NULL,
  created_at  timestamptz DEFAULT now(),
  UNIQUE (message_id, user_id, emoji)
);

CREATE TABLE activity_likes (
  activity_id uuid REFERENCES activities(id) ON DELETE CASCADE,
  user_id     uuid REFERENCES users(id) ON DELETE CASCADE,
  created_at  timestamptz DEFAULT now(),
  PRIMARY KEY (activity_id, user_id)
);

CREATE TABLE notifications (
  id          uuid PRIMARY KEY DEFAULT gen_random_uuid(),
  user_id     uuid REFERENCES users(id) ON DELETE CASCADE,
  type        text,  -- ACTIVITY_REQUEST|APPROVED|REJECTED|NEW_MESSAGE
  title       text NOT NULL,
  body        text,
  entity_type text,
  entity_id   uuid,
  is_read     boolean DEFAULT false,
  created_at  timestamptz DEFAULT now()
);
\end{lstlisting}

\textbf{Realtime:} Tin nhắn group chat được phân phối qua \textbf{WebSocket
Gateway} (NestJS) và \textbf{Redis Pub/Sub}. Mỗi activity có một Redis channel
riêng theo \texttt{activity:\{id\}}. Khi thành viên được duyệt, server tự động
subscribe user vào channel tương ứng.

\begin{lstlisting}[language={}]
-- Phan phoi tin nhan group chat
Channel: activity:{activity_id}
Payload: { userId, username, message, timestamp }

-- Server subscribe user sau khi duyet
redis.subscribe(`activity:${activityId}`)
\end{lstlisting}

% -------------------------------------------------------------------------
\subsection{Mạng xã hội Chia sẻ Lịch trình (Shared Trips)}

\subsubsection{Bài toán}

Sau khi tạo lập được một lịch trình tối ưu bằng công cụ Trip Planner, khách du
lịch thường có nhu cầu chia sẻ thành quả này lên cộng đồng để các du khách
khác tham khảo hoặc tái sử dụng. Các nền tảng hiện tại (như Facebook hay các
blog du lịch) chỉ hỗ trợ chia sẻ dưới dạng văn bản và hình ảnh tĩnh, khiến
người đọc muốn đi theo phải tự tay sao chép và thiết lập lại lịch trình từ đầu
rất mất thời gian.

Hệ thống cung cấp một cơ chế mạng xã hội du lịch thu nhỏ, liên kết trực
tiếp giữa bài viết và dữ liệu bản đồ/lịch trình thực tế.

\subsubsection{Luồng hoạt động và Tính năng}

\begin{enumerate}
  \item \textbf{Đăng bài chia sẻ lịch trình:} Trong giao diện quản lý Trip,
        người dùng chọn lịch trình đã lên và bấm chia sẻ lên cộng đồng. Hệ
        thống tạo một bản ghi \texttt{Activity} liên kết \texttt{trip\_id}
        tương ứng --- bài đăng chia sẻ và activity nhóm dùng chung một thực
        thể, phân biệt bởi sự hiện diện của \texttt{trip\_id}.
  \item \textbf{Phân luồng bảng tin (Feed Tabs):} Bảng tin hoạt động trên
        Frontend (\texttt{/activities}) được chia làm 3 tab chuyên biệt:
        \begin{itemize}
          \item \textbf{Groups:} Chỉ hiển thị các hoạt động rủ rê ghép nhóm
                du lịch thông thường (không đính kèm Trip). Các row hiển thị
                nút ``Join Group''.
          \item \textbf{Shared Trips:} Chỉ hiển thị các bài đăng chia sẻ
                lịch trình du lịch (có đính kèm \texttt{trip\_id}). Nút
                ``Join Group'' được thay thế bằng nút ``Save Trip''.
          \item \textbf{Joined:} Hiển thị danh sách các hoạt động hoặc lịch
                trình chia sẻ mà người dùng hiện tại đã tạo hoặc tham
                gia/lưu lại.
        \end{itemize}
  \item \textbf{Tương tác xã hội:}
        \begin{itemize}
          \item \textbf{Like (Yêu thích):} Cho phép người dùng thả tim bài
                đăng. Tính năng này áp dụng cơ chế \textbf{Optimistic UI}
                trên Frontend --- cập nhật tức thời trạng thái thích và đếm
                số lượt thích trên giao diện, đồng thời gửi request ngầm đến
                API. Nếu API báo lỗi, UI sẽ tự động phục hồi (rollback).
          \item \textbf{Comment (Bình luận):} Hệ thống tái sử dụng thực thể
                \texttt{Message} trên bảng \texttt{activities} để lưu trữ
                và hiển thị nội dung bình luận.
        \end{itemize}
  \item \textbf{Sao chép lịch trình (Cloning Itinerary):} Khi bấm nút ``Save
        Trip'', hệ thống kích hoạt API nhân bản lịch trình
        (\texttt{cloneActivityTrip}). Toàn bộ thông tin lịch trình gốc bao
        gồm tên (gắn hậu tố ``(Cloned)''), số ngày, điểm khởi hành, danh
        sách \texttt{trip\_days} và \texttt{trip\_stops} được sao chép sang
        tài khoản của người dùng mới. Trường \texttt{saves\_count} trên bài
        đăng tăng thêm 1.
\end{enumerate}

\subsubsection{Cơ chế nhân bản lịch trình}

Luồng xử lý tại Backend sử dụng dịch vụ NestJS phối hợp với Prisma ORM đảm
bảo tính nhất quán của dữ liệu theo mô hình giao dịch sâu:

\begin{enumerate}
  \item \textbf{Xác thực thông tin:} Kiểm tra tính hợp lệ của bài đăng
        activity và truy vấn đầy đủ cấu trúc lịch trình gốc gắn liền.
  \item \textbf{Truy vấn sâu (Deep Query):} Đọc toàn bộ dữ liệu phân tầng
        bao gồm lịch trình chính, \texttt{trip\_days} và \texttt{trip\_stops}
        theo đúng thứ tự hành trình.
  \item \textbf{Nhân bản sâu (Deep Clone):} Tạo mới Trip cho người dùng
        hiện tại, tự động gán hậu tố ``(Cloned)'', thiết lập
        \texttt{clonedFromId} tham chiếu về lịch trình gốc.
  \item \textbf{Đồng bộ thống kê:} Cập nhật tăng \texttt{saves\_count} trên
        bài viết chia sẻ và hoàn tất transaction.
\end{enumerate}

% -------------------------------------------------------------------------
\subsection{AI Mentor}

\subsubsection{Bài toán}

Khách du lịch cần một trợ lý ảo am hiểu địa phương để gợi ý địa điểm, giải đáp
thắc mắc về văn hóa Hà Nội, và đặc biệt \textbf{lập lịch trình tự động qua
hội thoại} --- thay thế hoàn toàn luồng Trip Planner thủ công bằng cách thu
thập thông số qua chat theo kiểu Socratic.

\subsubsection{Triển khai}

Module sử dụng mô hình \textbf{Gemini 2.5 Flash} (Google AI SDK) nhờ khả năng
xử lý context lớn và độ trễ thấp. Kiến trúc \textbf{Single Mega Prompt} --- một
lời nhắc duy nhất xử lý tất cả ý định, trả về JSON có cấu trúc:

\begin{lstlisting}[language={}]
{
  "intent": "nearby | trip_plan | trip_plan_collecting | general",
  "response": "<chuoi Markdown tieng Viet, than thien, emoji>",
  "extractedParams": {
    "numDays":         <integer | null>,
    "travelDate":      <"YYYY-MM-DD" | null>,
    "startTime":       <so nguyen 24h | null>,
    "endTime":         <so nguyen 24h | null>,
    "lunchBreakStart": <so nguyen 24h | null>,
    "lunchBreakEnd":   <so nguyen 24h | null>,
    "visitDuration":   <phut, 15-120 | null>,
    "districts":       <string[] | []>,
    "categories":      <string[] | []>
  }
}
\end{lstlisting}

\textbf{Bốn intent được hỗ trợ:}

\begin{itemize}
  \item \textbf{\texttt{nearby}:} User hỏi về địa điểm gần vị trí hiện tại.
        Backend tính khoảng cách Haversine từ GPS user đến 65 địa điểm, trả
        về 12 địa điểm gần nhất kèm tên, category, địa chỉ, khoảng cách.
        AI phản hồi bằng gợi ý tự nhiên, đồng thời trả về \texttt{places}
        (pin map) để Frontend hiển thị trên bản đồ.
  \item \textbf{\texttt{trip\_plan}:} Đủ 3 thông số bắt buộc (numDays,
        travelDate, categories). Backend gọi thẳng \texttt{TripPlannerService}
        để sinh lịch trình, trả về \texttt{itinerary} đầy đủ.
  \item \textbf{\texttt{trip\_plan\_collecting}:} Thiếu ít nhất 1 thông số.
        AI hỏi \textbf{tất cả thông số còn thiếu trong một lượt duy nhất}
        (không hỏi từng cái một). Frontend lưu \texttt{accumulatedParams} và
        gửi lại ở lượt tiếp theo để AI biết đã thu thập được gì.
  \item \textbf{\texttt{general}:} Câu hỏi thông thường về Hà Nội, văn hoá,
        lịch sử, ẩm thực.
\end{itemize}

\textbf{Thông số tuỳ chọn bổ sung} (so với phiên bản cơ sở):

\begin{itemize}
  \item \textbf{\texttt{lunchBreakStart / lunchBreakEnd}:} Cho phép user chỉ
        định khung giờ nghỉ trưa, truyền thẳng vào Trip Planner thay cho
        giờ nghỉ mặc định toàn hệ thống.
  \item \textbf{\texttt{visitDuration}:} Số phút tham quan mỗi địa điểm
        (15--120 phút), cho phép tuỳ chỉnh thay vì dùng
        \texttt{visit\_duration\_min} mặc định của từng địa điểm.
  \item \textbf{\texttt{districts}:} Lọc địa điểm theo khu vực Hà Nội.
\end{itemize}

\begin{lstlisting}[language={}]
System Context (built dynamically per request):
  Today: YYYY-MM-DD
  Recent conversation (last 8 turns)
  Already collected params (to avoid re-asking)
  Nearby places list (pre-sorted by distance, if applicable)

Intent Rules:
  - "trip_plan"            -> all 3 required params present
  - "trip_plan_collecting" -> missing >= 1 required param
                             -> ask ALL missing in ONE message
  - "nearby"               -> user asks about places near them
  - "general"              -> everything else

Output: JSON only, no markdown fences
\end{lstlisting}

% =============================================================================
\section{Kiến trúc Hệ thống}
% =============================================================================

\subsection{Tổng quan}

\begin{center}
\begin{tabular}{ll}
  \toprule
  \textbf{Tầng} & \textbf{Công nghệ} \\
  \midrule
  Frontend         & Next.js 14, Tailwind CSS, Leaflet.js, Zustand \\
  Backend API      & NestJS (Node.js), WebSocket Gateway (Socket.io) \\
  ORM              & Prisma (PostgreSQL extensions: postgis)          \\
  AI Integration   & Gemini 2.5 Flash (Google AI SDK)                 \\
  Routing Engine   & Goong Distance Matrix API / OSRM (self-hosted)  \\
  Database Chính   & PostgreSQL + PostGIS                             \\
  Cache / Realtime & Redis Pub/Sub                                    \\
  Maps             & CartoDB Voyager + Leaflet                        \\
  Email            & Nodemailer (Gmail SMTP)                          \\
  Media Upload     & Multer (NestJS), lưu \texttt{public/uploads}     \\
  \bottomrule
\end{tabular}
\end{center}

\subsection{Cấu trúc Backend (NestJS Modules)}

\begin{center}
\begin{tabular}{ll}
  \toprule
  \textbf{Module} & \textbf{Chức năng} \\
  \midrule
  \texttt{auth}       & Đăng ký, OTP, đăng nhập, refresh, reset password \\
  \texttt{users}      & Đọc/cập nhật hồ sơ người dùng \\
  \texttt{admin}      & Dashboard, User/Place/Report management, AdminLog \\
  \texttt{places}     & CRUD địa điểm, tìm kiếm PostGIS \\
  \texttt{trips}      & Trip Planner 5-phase + lưu lịch trình \\
  \texttt{activities} & CRUD activity, join/approve, likes, comments, clone \\
  \texttt{group-chat} & WebSocket Gateway, Redis Pub/Sub \\
  \texttt{ai-chat}    & Gemini integration, multi-intent routing \\
  \texttt{media}      & Upload avatar và ảnh activity (Multer) \\
  \texttt{prisma}     & Prisma service (shared across modules) \\
  \texttt{common}     & Guards, decorators, filters dùng chung \\
  \bottomrule
\end{tabular}
\end{center}

\subsection{Cấu trúc Frontend (Next.js App Router)}

\begin{center}
\begin{tabular}{ll}
  \toprule
  \textbf{Route} & \textbf{Nội dung} \\
  \midrule
  \texttt{/}                   & Landing page \\
  \texttt{/(auth)/register}    & Đăng ký \\
  \texttt{/(auth)/verify-email} & Xác thực OTP \\
  \texttt{/(auth)/login}       & Đăng nhập \\
  \texttt{/(auth)/forgot-password} & Quên mật khẩu \\
  \texttt{/(auth)/reset-password}  & Đặt lại mật khẩu \\
  \texttt{/(main)/discovery}   & Khám phá địa điểm + bản đồ \\
  \texttt{/(main)/places/:id}  & Chi tiết địa điểm \\
  \texttt{/(main)/trips}       & Trip Planner + lịch trình đã lưu \\
  \texttt{/(main)/activities}  & Activity feed (Groups/Shared/Joined) \\
  \texttt{/(main)/profile}     & Hồ sơ cá nhân \\
  \texttt{/(main)/profile/edit} & Chỉnh sửa hồ sơ \\
  \texttt{/admin/dashboard}    & Admin dashboard stats \\
  \texttt{/admin/users}        & Quản lý tài khoản \\
  \texttt{/admin/places}       & Quản lý địa điểm \\
  \texttt{/admin/reports}      & Quản lý báo cáo vi phạm \\
  \bottomrule
\end{tabular}
\end{center}

\subsection{Luồng dữ liệu chính}

\begin{enumerate}
  \item \textbf{Khám phá địa điểm:} Frontend gửi filter → NestJS truy vấn
        PostgreSQL (Prisma) → trả về danh sách + GeoJSON để render Leaflet
        pins.

  \item \textbf{Trip Planner:} Frontend gửi danh sách địa điểm và số ngày →
        NestJS tiền lọc địa điểm theo \texttt{open\_days} (Bước~0) → K-Means
        clustering (Deterministic với Seeded PRNG) → gọi Goong Distance Matrix
        API lấy ma trận khoảng cách thực tế → GNN Routing có tích hợp hàm chi
        phí ưu tiên di chuyển và xử lý giờ nghỉ trưa
        ($Cost = TravelTime \times 2 + T_{wait} \times 60$) → chạy Conflict
        Resolution để điền vào các khe trống (Gap Insertion) → lưu lịch trình
        vào \texttt{trips}/\texttt{trip\_days}/\texttt{trip\_stops} → trả về
        lịch trình JSON kèm danh sách \texttt{unscheduled}.

  \item \textbf{Activity \& Group Chat:} User tạo activity (Goong Autocomplete
        + LocationPickerMap) → lưu DB + render pin trên map → user khác request
        tham gia → tạo \texttt{Notification} cho host → chủ duyệt → NestJS
        subscribe user vào Redis channel → WebSocket đồng bộ group chat
        realtime + Message Reactions.

  \item \textbf{Chia sẻ Lịch trình (Shared Trips):} User chọn chia sẻ lịch
        trình $\to$ hệ thống tạo \texttt{Activity} liên kết \texttt{trip\_id}
        $\to$ bài đăng xuất hiện trên tab Shared Trips $\to$ User khác tương
        tác (Like dùng Optimistic UI / Comment) hoặc bấm ``Save Trip'' $\to$
        NestJS thực thi nhân bản sâu lịch trình (\texttt{cloneActivityTrip})
        kèm \texttt{clonedFromId} và tăng \texttt{saves\_count}.

  \item \textbf{AI Mentor:} Frontend gửi message + GPS + \texttt{accumulatedParams}
        → NestJS xây dựng Single Mega Prompt → Gemini 2.5 Flash trả về JSON
        có \texttt{intent} → nếu \texttt{trip\_plan}: gọi
        \texttt{TripPlannerService} trực tiếp → trả về \texttt{itinerary} +
        \texttt{accumulatedParams} mới cho client lưu lại.

  \item \textbf{Admin:} Admin đăng nhập tại \texttt{/admin/login} → JWT chứa
        \texttt{role=ADMIN} → truy cập Dashboard/Users/Places/Reports → mọi
        thao tác (ban, resolve report, delete) đều ghi \texttt{AdminLog} và
        thực hiện trong Prisma transaction.
\end{enumerate}

% =============================================================================
\section{Dữ liệu \& Nội dung}
% =============================================================================

Dự án sử dụng \textbf{một bộ dữ liệu duy nhất} (65 địa điểm du lịch trọng điểm)
đã được xác thực tọa độ và thông tin, phục vụ cho cả ba module chính, tránh phân
tán công sức thu thập và bảo trì. Dữ liệu được tổng hợp từ:

\begin{itemize}
  \item \textbf{OpenStreetMap (Overpass API)} --- toạ độ, tên, phân loại
        (giấy phép ODbL, miễn phí). Dữ liệu OSM Việt Nam cũng dùng để khởi
        tạo OSRM routing engine.
  \item \textbf{Bổ sung thủ công} --- giờ mở cửa, thời gian tham quan ước
        tính, ảnh đại diện, mô tả ngắn bằng tiếng Anh, tags phân loại.
\end{itemize}

\textbf{Phạm vi dữ liệu:} Chỉ bao gồm địa điểm du lịch (di tích, bảo tàng,
hồ, công viên, chùa,...). Hotel và nhà hàng không nằm trong dataset tĩnh ---
được tra cứu realtime bởi AI Mentor khi cần.

\medskip
Phân loại theo \textbf{khu vực} phục vụ clustering và filter:
Hoàn Kiếm, Ba Đình, Tây Hồ, Đống Đa, Hai Bà Trưng, Cầu Giấy.

\medskip
Phân loại theo \textbf{category} (dùng nhất quán xuyên suốt cả Discovery
filter, Activity category, và AI Mentor intent extraction):
Heritage \& History, Arts \& Culture, Sightseeing, Eat \& Shop,
Nature \& Outdoors, Spiritual.

\end{document}